% Generated draft: requires mathematical and source review.
\documentclass[11pt]{article}
\usepackage[margin=25mm]{geometry}
\usepackage{fontspec}
\setmainfont{Latin Modern Roman}
\usepackage{amsmath,amssymb,mathtools,unicode-math}
\setmathfont{Latin Modern Math}
\usepackage{enumitem,needspace,xurl,hyperref}
\setlength{\parindent}{0pt}
\setlength{\parskip}{5pt}
\setlength{\emergencystretch}{3em}
\newcommand{\N}{\mathbb{N}}
\newcommand{\Z}{\mathbb{Z}}
\newcommand{\Q}{\mathbb{Q}}
\newcommand{\R}{\mathbb{R}}
\newcommand{\C}{\mathbb{C}}
\newcommand{\F}{\mathbb{F}}
\newcommand{\A}{\mathbb{A}}
\newcommand{\PP}{\mathbb{P}}
\newcommand{\E}{\mathbb{E}}
\newcommand{\eps}{\varepsilon}
\newcommand{\dd}{\,\mathrm{d}}
\begin{document}
\section*{UnsolvedMath ID: 800772}
\textbf{Source problem code:} catalogue entry
\subsection{Definitions and mathematical statement}
Let \(X\) be a smooth projective variety, and let \(T_X=\mathcal{H}om_{\mathcal O_X}(\Omega_X^1,\mathcal O_X)\) denote its tangent bundle. The supplied statement assumes that \(T_X\) is pseudo-effective, but does not specify the definition of pseudo-effectivity for vector bundles or the ground field.

Under the usual characteristic-zero interpretation, a maximal rationally connected (MRC) fibration is a dominant rational map \(\phi\colon X\dashrightarrow Y\) whose general fibers are the equivalence classes of general points under connection by chains of rational curves. Its base is determined only up to birational equivalence. Thus the source's notation \(\phi\colon X\to Y\) requires clarification as to whether a morphism on all of \(X\) is assumed or is part of the requested conclusion.

For a specified topology and a morphism \(f\colon X\to Y\), local triviality means that there is a fiber \(F\) such that every point \(y\in Y\) has a neighborhood \(U\) and an isomorphism
\[
\theta_U\colon f^{-1}(U)\xrightarrow{\sim}U\times F,
\qquad \operatorname{pr}_U\circ\theta_U=f|_{f^{-1}(U)}.
\]
For the \(\acute{e}\)tale topology, neighborhoods are replaced by an \(\acute{e}\)tale covering and inverse images by base changes. In the complex analytic interpretation, the displayed isomorphisms are biholomorphic.

\textbf{Question (Matsumura, Problem 7, 2022).} If \(X\) is a smooth projective variety with pseudo-effective tangent bundle, is its MRC fibration \(\phi\colon X\to Y\) locally trivial?

The question is retained with unresolved qualifications: the ground field, the intended vector-bundle positivity convention, the topology of local triviality, the choice of MRC model, and whether the conclusion includes existence of an everywhere-defined morphism must be fixed before this becomes a fully specified assertion. Local triviality over the entire base and local triviality over a dense open subset are distinct possible targets; the supplied excerpt does not resolve this distinction.
\subsection{Short English statement}
Matsumura asks whether a smooth projective variety whose tangent bundle is pseudo-effective has a locally trivial maximal rationally connected fibration. The supplied excerpt leaves the ground field, positivity convention, topology, and treatment of the possibly rational fibration unspecified. This is a statement-only draft requiring mathematical review.
\subsection{Sources}
\begin{itemize}
\item Local UnsolvedMath registry, permanent ID \(800772\), source problem code \texttt{catalogue entry}, title: \emph{Positively Curved Projective Varieties --- Problem 7}. The supplied record attributes the question to Matsumura (2022) and labels it \texttt{open}; that label is supplied registry metadata, not an independently verified mathematical or literature assessment.
\item Matsumura, \emph{Open Problems on Structure of Positively Curved Projective Varieties} (2022), Problem 7. Supplied source URL: \url{https://arxiv.org/abs/2211.09153}. The record reports extraction by \texttt{arxiv\_latex} and access on \texttt{2026-10-16}. This access date is supplied provenance; no independent source visit or verification is claimed.
\end{itemize}
\end{document}
